\section{Методы и инструментарий}

\subsection{Спектральный анализ}
Устойчивость периодических орбит исследуется с помощью матрицы монодромии \(M\), которая получается интегрированием уравнений в вариациях за один период \(T\). Собственные значения \(\mu\) матрицы \(M\) называются мультипликаторами Флоке. В гамильтоновых системах они обладают свойством симметрии: если \(\mu\) — мультипликатор, то \(\mu^{-1}\), \(\bar{\mu}\) и \(\bar{\mu}^{-1}\) также являются мультипликаторами.

Все орбиты в каталогах классифицируются по спектру на три типа:
\begin{itemize}
    \item \textbf{C} — четыре тривиальных мультипликатора \(\mu \approx 1\) и две комплексно-сопряжённые пары вне единичной окружности;
    \item \textbf{R} — четыре тривиальных мультипликатора и четыре вещественных мультипликатора, образующих две пары \((\mu, 1/\mu)\);
    \item \textbf{OTHER} — четыре тривиальных мультипликатора, одна комплексная пара и одна вещественная пара.
\end{itemize}
Показатель Ляпунова определяется как \(\lambda = \frac{1}{T} \ln(\max |\mu|)\).

\subsection{Пространство форм и поверхность нулевой скорости}
Для визуализации и анализа топологии относительного движения вводится координатное пространство \(\xi = (\xi_1, \xi_2, \xi_3)\), основанное на векторах Якоби. Пусть \(Q_1 = r_2 - r_1\) и \(Q_2 = r_3 - \frac{m_1 r_1 + m_2 r_2}{m_1 + m_2}\). Тогда координаты определяются как:
\begin{equation}
    \xi_1 = \mu_1 |Q_1|^2 - \mu_2 |Q_2|^2, \quad
    \xi_2 + i\xi_3 = 2\sqrt{\mu_1 \mu_2} \, Q_1 \overline{Q_2}.
\end{equation}
где \(\mu_1 = \frac{m_1 m_2}{m_1+m_2}\), \(\mu_2 = \frac{m_3(m_1+m_2)}{m_1+m_2+m_3}\). В этих координатах точка равновесия имеет вид \(\xi_{eq} = (-1, \sqrt{3}, 0)\), а момент инерции выражается как \(I = |\xi| = \sqrt{\xi_1^2 + \xi_2^2 + \xi_3^2}\).

Поверхность нулевой скорости (ZVS) в координатах \(\xi\) задаётся уравнением Сундмана:
\begin{equation}
    V(\xi) + \frac{L_z^2}{2|\xi|} = h_{in},
\end{equation}
где \(V(\xi)\) — потенциал, \(L_z\) — полный момент импульса, \(h_{in}\) — полная энергия в инерциальной системе. Область, где левая часть меньше \(h_{in}\), является доступной для движения.
